\documentclass[10pt,a4paper]{amsart}
\usepackage{amsmath,amssymb,amsthm,graphicx}
% Retained minimal preamble for the Stacks Project chapter more-algebra.tex (More on Algebra).
% The chapter file pinned-original-source.tex carries only \input{preamble} on its first line, so no
% preamble text lives inside the chapter itself.  The Stacks Project uses one single shared preamble for
% every chapter; the snapshot stages it identically as sources/stacks-chow/preamble.tex and
% sources/stacks-etale-cohomology/preamble.tex.  The block below is copied mechanically, line for line,
% from sources/stacks-chow/preamble.tex lines 14-18 (the Xy-pic setup) and 154-217 (the native theorem
% declarations together with the mathematical macros).  Nothing was added, removed or redefined; the
% wrapper class and the amsmath/amssymb/amsthm/graphicx base packages are supplied by the assembler.
% Chapter: more-algebra.tex -- More on Algebra.
% For commutative diagrams we use Xy-pic
\usepackage[all]{xy}

% We use 2cell for 2-commutative diagrams.
\xyoption{2cell}

% Theorem environments.
%
\theoremstyle{plain}
\newtheorem{theorem}[subsection]{Theorem}
\newtheorem{proposition}[subsection]{Proposition}
\newtheorem{lemma}[subsection]{Lemma}

\theoremstyle{definition}
\newtheorem{definition}[subsection]{Definition}
\newtheorem{example}[subsection]{Example}
\newtheorem{exercise}[subsection]{Exercise}
\newtheorem{situation}[subsection]{Situation}

\theoremstyle{remark}
\newtheorem{remark}[subsection]{Remark}
\newtheorem{remarks}[subsection]{Remarks}

\numberwithin{equation}{subsection}

% Macros
%
\def\lim{\mathop{\mathrm{lim}}\nolimits}
\def\colim{\mathop{\mathrm{colim}}\nolimits}
\def\Spec{\mathop{\mathrm{Spec}}}
\def\Hom{\mathop{\mathrm{Hom}}\nolimits}
\def\Ext{\mathop{\mathrm{Ext}}\nolimits}
\def\SheafHom{\mathop{\mathcal{H}\!\mathit{om}}\nolimits}
\def\SheafExt{\mathop{\mathcal{E}\!\mathit{xt}}\nolimits}
\def\Sch{\mathit{Sch}}
\def\Mor{\mathop{\mathrm{Mor}}\nolimits}
\def\Ob{\mathop{\mathrm{Ob}}\nolimits}
\def\Sh{\mathop{\mathit{Sh}}\nolimits}
\def\NL{\mathop{N\!L}\nolimits}
\def\CH{\mathop{\mathrm{CH}}\nolimits}
\def\proetale{{pro\text{-}\acute{e}tale}}
\def\etale{{\acute{e}tale}}
\def\QCoh{\mathit{QCoh}}
\def\Ker{\mathop{\mathrm{Ker}}}
\def\Im{\mathop{\mathrm{Im}}}
\def\Coker{\mathop{\mathrm{Coker}}}
\def\Coim{\mathop{\mathrm{Coim}}}

% Boxtimes
%
\DeclareMathSymbol{\boxtimes}{\mathbin}{AMSa}{"02}

%
% Macros for moduli stacks/spaces
%
\def\QCohstack{\mathcal{QC}\!\mathit{oh}}
\def\Cohstack{\mathcal{C}\!\mathit{oh}}
\def\Spacesstack{\mathcal{S}\!\mathit{paces}}
\def\Quotfunctor{\mathrm{Quot}}
\def\Hilbfunctor{\mathrm{Hilb}}
\def\Curvesstack{\mathcal{C}\!\mathit{urves}}
\def\Polarizedstack{\mathcal{P}\!\mathit{olarized}}
\def\Complexesstack{\mathcal{C}\!\mathit{omplexes}}
% \Pic is the operator that assigns to X its picard group, usage \Pic(X)
% \Picardstack_{X/B} denotes the Picard stack of X over B
% \Picardfunctor_{X/B} denotes the Picard functor of X over B
\def\Pic{\mathop{\mathrm{Pic}}\nolimits}
\def\Picardstack{\mathcal{P}\!\mathit{ic}}
\def\Picardfunctor{\mathrm{Pic}}
\def\Deformationcategory{\mathcal{D}\!\mathit{ef}}

% --- appended 2026-10-12 (Chief): the project own definitions of the reference/slogan/history
% display environments, copied mechanically from the shared Stacks preamble lines 8-12
% (\usepackage{verbatim} plus the three \newenvironment lines, which the chapter preamble keeps
% verbatim from the project preamble). Nothing added, removed or redefined: these environments
% carry no mathematical content and the project itself typesets them as comments. Without these
% three lines a counted statement whose own text uses them fails with Environment slogan undefined.
\usepackage{verbatim}
\newenvironment{reference}{\comment}{\endcomment}
\newenvironment{slogan}{\comment}{\endcomment}
\newenvironment{history}{\comment}{\endcomment}
\usepackage{cleveref}
\crefname{originallemma}{original lemma}{original lemmas}
\Crefname{originallemma}{Original lemma}{Original lemmas}
\makeatletter
\newlabel{algebra-lemma-long-exact-sequence-tor}{{00M0}{source}{}{sourcepointer-1}{}}
\newlabel{algebra-lemma-long-exact-sequence-tor@cref}{{[originallemma][0][]00M0}{[1][0][]source}{}{}{}}
\makeatother
\crefname{originallemma}{original lemma}{original lemmas}
\Crefname{originallemma}{Original lemma}{Original lemmas}
\makeatletter
\newlabel{algebra-lemma-tor-commutes-filtered-colimits}{{0BNF}{source}{}{sourcepointer-2}{}}
\newlabel{algebra-lemma-tor-commutes-filtered-colimits@cref}{{[originallemma][0][]0BNF}{[1][0][]source}{}{}{}}
\makeatother
\crefname{originallemma}{original lemma}{original lemmas}
\Crefname{originallemma}{Original lemma}{Original lemmas}
\makeatletter
\newlabel{lemma-first-tor-total}{{25486}{source}{}{sourcepointer-3}{}}
\newlabel{lemma-first-tor-total@cref}{{[originallemma][0][]25486}{[1][0][]source}{}{}{}}
\makeatother
\crefname{originallemma}{original lemma}{original lemmas}
\Crefname{originallemma}{Original lemma}{Original lemmas}
\makeatletter
\newlabel{lemma-same-f-torsion}{{25255}{source}{}{sourcepointer-4}{}}
\newlabel{lemma-same-f-torsion@cref}{{[originallemma][0][]25255}{[1][0][]source}{}{}{}}
\makeatother
\crefname{originallemma}{original lemma}{original lemmas}
\Crefname{originallemma}{Original lemma}{Original lemmas}
\makeatletter
\newlabel{lemma-torsion-completion}{{24173}{source}{}{sourcepointer-5}{}}
\newlabel{lemma-torsion-completion@cref}{{[originallemma][0][]24173}{[1][0][]source}{}{}{}}
\makeatother
\title{For a glueing pair the first Tor of the bigger ring with the cokernel of localizing a module at the element vanishes}
\author{The Stacks Project Authors}
\date{stacks-project:more-algebra.tex:local-snapshot-20261012; GNU FDL 1.2 or later}
\begin{document}
\maketitle
\setcounter{subsection}{1}
\noindent\textit{Editorial note.} All delivered mathematical text is the author's own, copied from the pinned Stacks Project chapter file more-algebra.tex; no statement or proof step was rewritten, reordered or removed.\par\medskip
\noindent\textit{Editorial note.} The chapter file pinned-original-source.tex carries only \textbackslash{}input\{preamble\} on its first line, so no preamble text lives inside the chapter. The retained file sources/stacks-more-algebra/retained-preamble.tex is a mechanical, line-for-line copy of the project's single shared preamble as staged in the snapshot, together with the project's own definitions of the reference/slogan/history display environments, which the project itself typesets as comments; nothing was added, removed or redefined by this card and no package was added.\par\medskip
\noindent\textit{Editorial note.} The source URL is the source-level Stacks Project URL, as required when more than one card is produced from one source; the per-claim tag is carried by claim\_key.\par\medskip
\noindent\textit{Editorial note.} The counted statement span includes the author's own \textbackslash{}begin\{reference\} block recording that the lemma is a slight generalization of a cited result; that block is delivered verbatim and produces no unresolved citation. The author writes 'we deduce  that' with a double space in the last paragraph, reproduced unchanged. The five references inside the counted proof span are recorded as non-load-bearing navigation pointers.\par\medskip
\noindent\textit{Allowed standard prerequisites (editorial).}\begin{itemize}
\item Glueing pairs (R -> R', f) and the modules R/R[f\textasciicircum{}infty] and R'/R'[f\textasciicircum{}infty]. the setting at statement line 25503 and used at proof lines 25518-25520. Fixes the objects appearing in the special case and the comparison M tensor\_R R' = R'/R'[f\textasciicircum{}infty].
\item Commutation of Tor with filtered colimits. invoked at proof lines 25513-25514. Reduces the statement for M\_f/M to the statements for the quotients M/f\textasciicircum{}nM.
\item The long exact sequence for Tor and the injectivity of the multiplication map in the displayed diagram. used at proof lines 25521-25536. Is the computation that kills Tor\_1 in the special case.
\item Exactness of tensoring with R' for the quotients involved. invoked at proof lines 25545-25546. Lets the general case be reduced to the free case.
\end{itemize}\medskip
\ifdefined\hypertarget\hypertarget{sourcepointer-1}{}\fi
\noindent\textit{Source reference.} Original lemma 00M0: focused-seeds/AGP/sources/stacks-snapshot-20261011/tags.txt, source lines 806--806. Cross-chapter navigation pointer into the Stacks Project; the standard fact it supplies is applied inside the counted proof span and the argument resting on it stays there.\par\medskip
\ifdefined\hypertarget\hypertarget{sourcepointer-2}{}\fi
\noindent\textit{Source reference.} Original lemma 0BNF: focused-seeds/AGP/sources/stacks-snapshot-20261011/tags.txt, source lines 14663--14663. Cross-chapter navigation pointer into the Stacks Project; the standard fact it supplies is applied inside the counted proof span and the argument resting on it stays there.\par\medskip
\ifdefined\hypertarget\hypertarget{sourcepointer-3}{}\fi
\noindent\textit{Source reference.} Original lemma 25486: sources/stacks-more-algebra/pinned-original-source.tex, source lines 25486--25486. Same-chapter navigation pointer; the statement it names is standard for the section and the argument that consumes it stays inside the counted proof span.\par\medskip
\ifdefined\hypertarget\hypertarget{sourcepointer-4}{}\fi
\noindent\textit{Source reference.} Original lemma 25255: sources/stacks-more-algebra/pinned-original-source.tex, source lines 25255--25255. Same-chapter navigation pointer; the statement it names is standard for the section and the argument that consumes it stays inside the counted proof span.\par\medskip
\ifdefined\hypertarget\hypertarget{sourcepointer-5}{}\fi
\noindent\textit{Source reference.} Original lemma 24173: sources/stacks-more-algebra/pinned-original-source.tex, source lines 24173--24173. Same-chapter navigation pointer; the statement it names is standard for the section and the argument that consumes it stays inside the counted proof span.\par\medskip
\begin{lemma}
\label{lemma-BL3}
\begin{reference}
Slight generalization of \cite[Lemme 3(a)]{Beauville-Laszlo}
\end{reference}
Let $(R \to R', f)$ be a glueing pair. For every $R$-module $M$, we have
$\text{Tor}^R_1(R', \Coker(M \to M_f)) = 0$.
\end{lemma}

\begin{proof}
Set $\overline{M} = M/M[f^\infty]$. Then
$\Coker(M \to M_f) \cong \Coker(\overline{M} \to \overline{M}_f)$
hence we may and do assume that $f$ is a nonzerodivisor on $M$.
In this case $M \subset M_f$ and $M_f/M = \colim M/f^nM$ where the
transition maps are given by multiplication by $f$. Since
formation of Tor groups commutes with colimits
(Algebra, Lemma \ref{algebra-lemma-tor-commutes-filtered-colimits})
it suffices to show that $\text{Tor}^R_1(R', M/f^n M) = 0$.

\medskip\noindent
We first treat the case $M = R/R[f^\infty]$. By
Lemma \ref{lemma-same-f-torsion}
we have $M \otimes_R R' = R'/R'[f^\infty]$.
From the short exact sequence $0 \to M \to M \to M/f^nM \to 0$
we obtain the exact sequence
$$
\xymatrix{
\text{Tor}_1^R(R', R/R[f^\infty]) \ar[r] &
\text{Tor}_1^R(R', M/f^n M) \ar[r] &
R'/R'[f^\infty] \ar[dll]_{f^n} \\
R'/R'[f^\infty] \ar[r] &
(R'/R'[f^\infty])/(f^n
(R'/R'[f^\infty])) \ar[r] & 0
}
$$
by Algebra, Lemma \ref{algebra-lemma-long-exact-sequence-tor}.
Here the diagonal arrow is injective. Since the first group
$\text{Tor}_1^R(R', R/R[f^\infty])$ is zero by
Lemma \ref{lemma-first-tor-total}, we deduce that
$\text{Tor}_1^R(R', M/f^nM) = 0$ as desired.

\medskip\noindent
To treat the general case, choose a surjection $F \to M$ with $F$ a free
$R/R[f^\infty]$-module, and form an exact sequence
$$
0 \to N \to F/f^n F \to M/f^n M \to 0.
$$
By Lemma \ref{lemma-torsion-completion}
this sequence remains unchanged, and hence
exact, upon tensoring with $R'$.
Since $\text{Tor}^R_1(R', F/f^n F) = 0$ by the
previous paragraph, we deduce  that
$\text{Tor}^R_1(R', M/f^n M) = 0$ as desired.
\end{proof}

\end{document}
